% ============================================================================
% 第 7 章 总结与展望
% ============================================================================
\section{总结与展望}
\label{sec:conclusion}

\subsection{全文工作总结}

本文面向自动化立体仓库中多取货任务并发下的多穿梭无人车（AGV）智能调度问题，围绕“任务指派”与“离散时空无碰撞路径规划”的耦合挑战，完成了从多层/混合货架形式化建模、核心算法演进、零依赖工程架构落地到大规模多随机种子实证检验的完整研究闭环。主要成果归纳如下：

\begin{enumerate}
  \item \textbf{构建了严格可验的三层协同架构与零依赖通用仿真平台（\code{src/agv/} v0.6.0）}：渐进建立了涵盖 Legacy 规则栅格、V4 单面多层货架（层高相关装货耗时与同面层序解锁）及 V9/V9.1 混合背靠背货架（显式选侧、开放流显式关闭与安全待命）的数学模型；底层通过增量索引式时空预约表、安全区间路径规划（SIPP）与线程局部几何缓存大幅压缩长时域避碰搜索开销；顶层实现解耦的独立物理验解器与本地交互式 Web 可视化界面（\code{agv serve}），确保所有实验输出 100\% 物理合法且支持一键复现。
  \item \textbf{针对多工况特征发展了 LNS 联合修复、BCRD 反事实影子试派与 P90 尾部保护机制}：在通用引擎主线 A 中，静态首轮批滚动（\code{sipp\_batch}）将 Legacy 静态流程时间改善 $+2.63\%$、总步数减少 $4.24\%$；高积压反事实影子试派（\code{bcrd\_budget}）在 V4 动态 5-seed 实验中取得 $+3.01\%$ 综合改善（单次耗时 $0.17\,\text{s}$），配合无损事件驱动影子推进在 90 个留出场景上削减 $43.2\%$ CPU 开销；在 V9.1 的 960 次未见种子留出实验中，历史热度安全选点与 ETA 匹配（\code{G4}）及流式 BCRD（\code{B0}/\code{B1}）分别在动态突发与密集积压工况下取得显著正收益，其中 P90 尾部护栏将留出集尾部退化例数从 13/240 清零至 0/240。
  \item \textbf{在 WEPAStacks 真实工业订单流上通过跨窗口多 Seed 复评确立了最终方案 \code{gate}}：在 $62\times 47$ 网格、6 车、5 层、900\,s 真实订单窗口及固定尺度线性目标下，经由 200 次开发窗与 156 次既有验证窗多随机种子配对实验证实，“事件门控滚动派单＋软分区联合 GA＋预约表 SIPP 物理预演”（\code{gate}）在全部 4 个既有验证窗口上稳定战胜 \code{fcfs}（$+5.45\%$）、\code{eta}（$+2.43\%$）与 \code{hungarian}（$+3.87\%$），并通过事件门控将开发窗 GA 触发次数削减 $63\%$；轻量档 \code{gate\_preview3} 可在 $F$ 仅变化 $-0.33\%$ 的条件下节省约 $44\%$ CPU 时间。同时，跨窗复评定量揭示了匈牙利初始注入（\code{hga}）、单开交叉自适应（\code{wide\_adapt01}）、变异自适应饱和（$>86\%$ 代数触顶 $0.95$）及过度扩大代数（\code{gate\_gen24}）在未知订单窗口上的过拟合回退机理，并完成了 Double DQN 元调度切换的对比归档。
\end{enumerate}

\begin{teammatebox}{7.1 节补充留白（可选 —— 供队友补充总结）}
  队友在完成 4.3、4.4、5.3、6.2 及 6.4.4 节内容后，可在此处补充关于 WEPA 真实流建模、\code{gate} 软分区编码设计及早期阶段消融实验的总结凝练（定稿时可直接合并入上方总结列表并移除本提示框）。
\end{teammatebox}

\subsection{局限性分析与未来展望}

尽管本系统在多类合成与真实工业订单窗口上展现出稳健的调度质量与工程可复现性，后续研究仍可沿以下方向进一步深化：
\begin{itemize}
  \item \textbf{更大规模车队与高拥堵窗口联合重排（RHCR 融合）}：当前实验主要聚焦 4--6 辆 AGV 与 3--5 层立体货架；当车队规模扩展至 12--20 辆且巷道密度进一步提高时，可结合滚动时域冲突消解（RHCR）~\cite{li2021rhcr} 在有限时间窗内引入周期性多车路径重规划。
  \item \textbf{残余尾估计引导的浅层影子试派}：BCRD 消融实验表明，若无可靠的残余完工时间估计，单纯截断影子推演深度易触发安全回退；未来可结合离线学习的价值函数拟合残余积压代价，在超大订单池下进一步压缩前瞻步数。
  \item \textbf{电量约束与多工作站动态协同}：在实际全天候仓储运行中，可进一步纳入 AGV 电池能耗模型、充电桩排队调度以及多拣选工作站动态负载均衡，推动算法向工业级全要素数字孪生系统演进。
\end{itemize}
